% Section: Experiments (claims + evidence; protocol in Appendix)
\section{Experiments}
\label{sec:experiments}

\subsection{Experimental Setup}
\label{sec:experiments-setup}

We evaluate AAR in two complementary settings.
The \emph{routing benchmark} is the primary development testbed: every query has offline
oracle outcomes for all four strategies, enabling regret-based comparison to the utility
oracle and controlled ablations of features and supervision.
\emph{Benchmark transfer} assesses whether a router trained on the routing benchmark deploys
well on full public workloads where only pre-execution features are available at inference and
strategies are executed live.
\textbf{We next ask} three questions:
(Q1)~Does cost-soft supervision reduce regret relative to hard labels?
(Q2)~Does fusing planner-derived complexity with semantic embeddings improve routing
beyond either signal alone (including a complexity-only ablation)?
(Q3)~Do routing-benchmark gains transfer to executed workloads without oracle access?

\paragraph{Benchmarks.}
We use two evaluation settings.
The \emph{routing benchmark} aggregates queries from multiple reasoning domains into
train, validation, and test partitions (split sizes in Appendix~\ref{app:benchmark-construction}).
For each training and evaluation query on this benchmark, all four strategies are executed
once offline to obtain correctness and cost labels; the deployed router never accesses
per-query oracle outcomes at inference.
We then assess \emph{benchmark transfer} under executed inference on five public workloads:
\textbf{GAIA}~\citep{mialon2023gaia}, \textbf{HotpotQA}~\citep{yang2018hotpotqa},
\textbf{MuSiQue}~\citep{trivedi2022musique}, \textbf{MMLU-Pro}~\citep{wang2024mmlupro},
and \textbf{MATH}~\citep{hendrycks2021math} (MMLU-Pro denoted MMLU in tables where space is limited).
Dataset selection rationale and workload descriptions are in
Appendix~\ref{app:benchmarks} (Tables~\ref{tab:deployment-sizes}--\ref{tab:pooled-utility});
routing-benchmark construction is in Appendix~\ref{app:benchmark-construction}.

\paragraph{Baselines and compared methods.}
All policies route among four strategies
(Eq.~\ref{eq:strategy-set}): direct prompting, chain-of-thought reasoning, ReAct-style tool
use, and multi-agent coordination.
\textbf{Fixed-strategy baselines} always execute one strategy and characterize non-adaptive
deployment.
\textbf{Learned routers} train tabular classifiers on $\phi(q)$ under hard labels,
utility-softmax targets, or cost-soft targets $p^\star(\cdot\mid q)$.
\textbf{Heuristic routers} replace planner-derived $\mathbf{z}(q)$ with hand-crafted lexical
and query-level heuristic features (Appendix~\ref{app:features}).
\textbf{Oracle references} include the utility-maximizing strategy $a^\star(q)$
(Eq.~\ref{eq:hard-oracle}) and a per-query \emph{best-agent} policy (Appendix~\ref{app:benchmarks}).

\paragraph{Metrics and training.}
The router is trained by minimizing
$\mathrm{KL}\!\bigl(p^\star(\cdot\mid q)\,\|\,p_\theta(\cdot\mid\phi(q))\bigr)$
under cost-soft targets $p^\star$ (Eq.~\ref{eq:cost-soft}), with hard-label and utility-softmax
alternatives in ablations.
On the routing benchmark we report utility regret (Eq.~\ref{eq:regret}), agent and total regret,
and executed exact match; on transfer workloads we report executed exact match and mean cost.
Representation ablations compare embedding-only, heuristic, complexity-only ($\mathbf{z}$), and
complexity--semantic ($\mathbf{z}+\psi$) configurations (Appendix~\ref{app:features}).
Training details, hyperparameters, and bootstrap protocols are in
Appendices~\ref{app:hyperparameters} and~\ref{app:additional-analyses}.
Unless stated otherwise, routing-benchmark tables refer to the held-out test partition ($n{=}100$).

\subsection{Summary of Findings}
\label{sec:results-overview}

We organize empirical claims into three layers.
\textbf{Main result (routing benchmark).}
On the held-out routing benchmark, \emph{cost-soft supervision} provides the largest reduction
in utility regret (agent regret $0.361\to 0.192$; Table~\ref{tab:hard-vs-soft}).
Planner-derived complexity further improves routing when \emph{combined} with semantic
embeddings: the fused representation attains the lowest total regret ($0.198$), while
complexity-only features underperform embedding-only baselines
(Table~\ref{tab:feature-regret}).
\textbf{Secondary result (benchmark transfer).}
Routing-benchmark improvements do not fully translate into exact-match gains across all
benchmark families (48.1\% pooled exact match vs.\ 50.3\% for a heuristic baseline;
Table~\ref{tab:live-benchmark}).
\textbf{Oracle reference (Best agent).}
Rows labeled \textbf{Best agent} quantify headroom from per-query oracle selection, not a
deployable policy.

\subsection{Representation Analysis}
\label{sec:representation-analysis}
\label{sec:feature-ablation}
\label{sec:main-results}

\textbf{Research question.}
Does planner-derived procedural complexity improve routing when fused with semantic
embeddings, relative to either signal alone?

\begin{table}[t]
  \centering
  \footnotesize
  \setlength{\tabcolsep}{2pt}
  \caption{Routing benchmark: representation comparison (held-out test, $n{=}100$).
    All rows: HGBM trained with cost-soft $p^\star$.
    Total regret includes QCE planner cost for complexity-only and complexity--semantic rows
    (embedding-only and heuristic rows do not invoke the planner).}
  \label{tab:feature-regret}
  \begin{tabular}{@{}lrrr@{}}
    \toprule
    Representation & Total regret $\downarrow$ & Agent regret $\downarrow$ & EM $\uparrow$ \\
    \midrule
    Embedding-only ($\psi$) & 0.242 & 0.242 & 77\% \\
    Heuristic baseline & 0.240 & 0.240 & 77\% \\
    Complexity-only ($\mathbf{z}$) & 0.273 & 0.267 & 74\% \\
    \textbf{Complexity--semantic ($\mathbf{z}+\psi$)} & \textbf{0.198} & \textbf{0.192} & \textbf{82\%} \\
    \bottomrule
  \end{tabular}
\end{table}

\begin{table}[t]
  \centering
  \footnotesize
  \caption{Alignment with oracle routing assignments (macro-F1, diagnostic).}
  \label{tab:feature-macro-f1}
  \begin{tabular}{@{}lr@{}}
    \toprule
    Representation & Macro-F1 $\uparrow$ \\
    \midrule
    Embedding-only ($\psi$) & 0.261 \\
    Heuristic baseline & 0.255 \\
    Complexity-only ($\mathbf{z}$) & 0.290 \\
    \textbf{Complexity--semantic ($\mathbf{z}+\psi$, proposed)} & \textbf{0.323} \\
    \bottomrule
  \end{tabular}
\end{table}

Table~\ref{tab:feature-regret} reports regret and executed exact match under cost-soft
supervision on the held-out routing benchmark; Table~\ref{tab:feature-macro-f1} reports
macro-F1 as a diagnostic of alignment with hard oracle routing assignments $a^\star(q)$
(not identical to utility regret under cost-soft routing).
Complexity-only improves macro-F1 over embedding-only and heuristics ($0.290$ vs.\
$0.261$/$0.255$) yet attains higher regret (Table~\ref{tab:feature-regret}); the proposed
fusion remains best on both metrics ($0.323$ macro-F1).
\textbf{These results suggest} that neither planner-derived complexity nor semantic embeddings
alone suffice for routing: complexity-only ($\mathbf{z}$) attains higher total and agent regret
than embedding-only (total $0.273$ vs.\ $0.242$; agent $0.267$ vs.\ $0.242$) and lower
executed exact match (74\% vs.\ 77\%).
Fusing $\mathbf{z}(q)$ with $\psi(q)$, however, yields the best configuration---total regret
$0.198$, agent regret $0.192$, and 82\% executed exact match---improving over embedding-only
by $0.044$ total / $0.050$ agent regret and over complexity-only by $0.075$ total /
$0.075$ agent regret.
\textbf{Takeaway.}
Complexity-only representations underperform embedding-only features, indicating that procedural
complexity is insufficient in isolation.
Planner-derived complexity is nonetheless \emph{complementary}: the full
$\phi(q)=[\mathbf{z}(q);\psi(q)]$ representation strictly outperforms either component alone
under our benchmark.

To examine which aspects of planner-derived complexity drive routing gains, we ablate the
five QCE dimensions $z_1,\ldots,z_5$ (Table~\ref{tab:complexity-dims}) while holding the
supervision protocol and semantic embedding $\psi(q)$ fixed.
Each ablated model uses \emph{one} complexity coordinate plus $\psi(q)$; the full system uses
all five coordinates with $\psi(q)$.
Table~\ref{tab:qce-dim-ablation} summarizes held-out test regret for each configuration;
Figure~\ref{fig:qce-dim-ablation} plots the single-dimension rows.

\begin{table}[t]
  \centering
  \footnotesize
  \caption{QCE dimension ablation (held-out test, $n{=}100$).
    HGBM with cost-soft $p^\star$; each row uses one complexity dimension plus $\psi(q)$.}
  \label{tab:qce-dim-ablation}
  \begin{tabular}{@{}lrr@{}}
    \toprule
    Feature & Total regret $\downarrow$ & Agent regret $\downarrow$ \\
    \midrule
    \textbf{Full (five dims + $\psi$)} & \textbf{0.198} & \textbf{0.192} \\
    Structural only ($z_1$) & 0.208 & 0.202 \\
    Reasoning only ($z_2$) & 0.210 & 0.204 \\
    Evidence only ($z_3$) & 0.228 & 0.222 \\
    Tool only ($z_4$) & 0.247 & 0.241 \\
    Coordination only ($z_5$) & 0.249 & 0.243 \\
    \bottomrule
  \end{tabular}
\end{table}

\begin{figure}[t]
  \centering
  \IfFileExists{results/reports/figures/fig6_qce_dimension_ablation.png}{%
    \includegraphics[width=0.92\linewidth]{results/reports/figures/fig6_qce_dimension_ablation.png}%
  }{%
    \fbox{\parbox[c][1.0in][c]{0.9\linewidth}{\centering\small
      Run \texttt{python -m research.figures} $\rightarrow$
      \texttt{fig6\_qce\_dimension\_ablation.png}.}}%
  }
  \caption{Single-dimension QCE ablation: total utility regret when routing with one
    complexity dimension plus semantic embeddings (dashed line: full five-dimensional
    $\mathbf{z}(q)$ with $\psi$).}
  \label{fig:qce-dim-ablation}
\end{figure}

Table~\ref{tab:qce-dim-ablation} and Figure~\ref{fig:qce-dim-ablation} show that no single
dimension recovers the performance of the full vector: total regret rises from $0.198$ to at
least $0.208$ when any one dimension is used alone.
Structural ($z_1$) and reasoning ($z_2$) dimensions yield the lowest single-signal regret
($0.208$ and $0.210$), consistent with routing dependence on anticipated plan depth and
inference demand.
Evidence ($z_3$) remains informative but weaker in isolation ($0.228$).
Tool ($z_4$) and coordination ($z_5$) dimensions are the least sufficient on their own
($0.247$ and $0.249$), suggesting that retrieval- and orchestration-related cues are most
valuable when combined with other procedural signals rather than as standalone routers.

Leave-one-out ablations in Appendix~\ref{app:qce-dim-loo} reinforce this picture: withholding
reasoning produces nearly the same regret as the full model ($0.199$), whereas withholding
structural or evidence coordinates increases regret more substantially ($0.228$ and $0.206$).
Together, the single-dimension and leave-one-out results indicate that QCE captures
\emph{complementary} facets of procedural structure; the five-dimensional design is justified
even though no single facet alone matches full fusion.

\subsection{Supervision Analysis}
\label{sec:supervision-analysis}
\label{sec:soft-target-ablation}
\label{sec:routing-benchmark}

\textbf{Research question.}
Does cost-soft supervision improve routing relative to hard labels and utility-softmax
targets?

\begin{table}[t]
  \centering
  \footnotesize
  \setlength{\tabcolsep}{3pt}
  \caption{Routing benchmark: hard vs.\ cost-soft supervision (held-out test).}
  \label{tab:hard-vs-soft}
  \begin{tabular}{@{}llrrr@{}}
    \toprule
    Policy target & Classifier & Total regret $\downarrow$ & Agent regret $\downarrow$ & EM $\uparrow$ \\
    \midrule
    Hard & HGBM & 0.366 & 0.361 & 63\% \\
    Hard & Logistic & 0.251 & 0.245 & 76\% \\
    Soft ($p^\star$) & Logistic & 0.205 & 0.199 & 81\% \\
    \textbf{Soft ($p^\star$)} & \textbf{HGBM} & \textbf{0.198} & \textbf{0.192} & \textbf{82\%} \\
    \bottomrule
  \end{tabular}
\end{table}

Table~\ref{tab:hard-vs-soft} compares hard-label and cost-soft training on the held-out test
split for the proposed representation.
\textbf{These results suggest} that supervision design dominates classifier capacity under our
setting: hard-label HGBM underperforms soft-label logistic regression, whereas cost-soft HGBM
achieves the best overall configuration.
Cost-soft supervision with gradient-boosted trees achieves the lowest regret, reducing agent
regret from $0.361$ to $0.192$ and raising executed exact match from 63\% to 82\%.
The supervision effect exceeds differences between classifiers under soft training,
indicating that \emph{how} the router is supervised matters more than the choice of
classifier alone.

\begin{table}[t]
  \centering
  \footnotesize
  \setlength{\tabcolsep}{3pt}
  \caption{Routing benchmark: supervision comparison (held-out test, $n{=}100$).
    All rows: HGBM, complexity--semantic $\phi(q)$.}
  \label{tab:soft-target}
  \begin{tabular}{@{}lrrr@{}}
    \toprule
    Target & Total regret $\downarrow$ & Agent regret $\downarrow$ & EM $\uparrow$ \\
    \midrule
    Hard labels (HGBM) & 0.366 & 0.361 & 63\% \\
    Utility-softmax & 0.243 & 0.237 & 77\% \\
    \textbf{Cost-soft $p^\star$} & \textbf{0.198} & \textbf{0.192} & \textbf{82\%} \\
    \bottomrule
  \end{tabular}
\end{table}

Table~\ref{tab:soft-target} contrasts supervision targets under the same representation.
Cost-soft $p^\star$ achieves the lowest regret ($0.198$ total, $0.192$ agent) and highest
executed exact match ($82\%$), whereas utility-softmax yields $0.243$/$0.237$ and hard
labels (HGBM) yield $0.366$/$0.361$.
\textbf{Takeaway.}
Restricting training mass to correct strategies and reweighting them by cost provides
better supervision than hard argmax labels or utility-softmax over all strategies.

\subsection{Benchmark Transfer Evaluation}
\label{sec:benchmark-transfer}
\label{sec:deployment-evaluation}

\textbf{Research question.}
Do routing policies learned on the routing benchmark transfer to executed workloads
without per-query oracle access at inference?

\begin{table*}[t]
  \centering
  \footnotesize
  \setlength{\tabcolsep}{4pt}
  \renewcommand{\arraystretch}{1.05}
  \caption{\textbf{Benchmark transfer evaluation.}
    Execution exact match (EM, \%) and mean cost (mUSD/query) across five executed benchmark families.
    Fixed-strategy baselines always execute one strategy.
    \textbf{Average} is the unweighted mean across fixed strategies.
    \textbf{Best agent} is an oracle reference that selects the lowest-cost correct strategy per query (or lowest-cost strategy if none is correct).
    \textbf{Selective} applies QCE-based gating.
    \textbf{Proposed}$^\dagger$ uses a complexity--semantic representation with cost-soft supervision;
    \textbf{Heuristic baseline} uses heuristic routing features.
    Costs include query-complexity estimation and strategy execution.}
  \label{tab:live-benchmark}
  \begin{tabular}{@{}l *{5}{cc}@{}}
    \toprule
    \multirow{2}{*}{Method} &
      \multicolumn{2}{c}{GAIA} &
      \multicolumn{2}{c}{MMLU} &
      \multicolumn{2}{c}{HotpotQA} &
      \multicolumn{2}{c}{MATH} &
      \multicolumn{2}{c}{MuSiQue} \\
    \cmidrule(lr){2-3}\cmidrule(lr){4-5}\cmidrule(lr){6-7}\cmidrule(lr){8-9}\cmidrule(lr){10-11}
    & EM & mUSD & EM & mUSD & EM & mUSD & EM & mUSD & EM & mUSD \\
    \midrule
    Raw        &  4.2 & 0.044 & 48.5 & 0.057 & 39.5 & 0.041 & 31.0 & 0.049 &  2.5 & 0.042 \\
    CoT        &  6.1 & 0.072 & 53.0 & 0.059 & 36.5 & 0.041 & 69.5 & 0.278 &  2.5 & 0.042 \\
    ReAct      & 26.1 & 4.709 & 59.0 & 0.134 & 69.0 & 1.133 & 63.0 & 0.658 & 36.5 & 2.486 \\
    Multiagent & 13.3 & 8.450 & 46.5 & 0.451 & 32.0 & 2.638 & 56.0 & 1.138 & 25.5 & 8.658 \\
    \midrule
    Average    & 12.4 & 3.319 & 51.8 & 0.175 & 44.2 & 0.963 & 54.9 & 0.531 & 16.8 & 2.807 \\
    Best agent & 31.5 & 1.063 & 70.5 & 0.093 & 76.5 & 0.419 & 78.0 & 0.180 & 46.0 & 1.343 \\
    Selective  & 18.8 & 3.241 & 54.5 & 0.278 & 61.5 & 1.019 & 67.5 & 0.548 & 36.0 & 3.279 \\
    \midrule
    Ours$^\dagger$ & 18.8 & 3.409 & 55.0 & 0.354 & 58.5 & 1.073 & 67.0 & 0.637 & 36.0 & 2.920 \\
    Heuristic  & 24.9 & 4.426 & 56.0 & 0.118 & 64.0 & 1.120 & 67.0 & 0.444 & 35.0 & 3.322 \\
    \bottomrule
  \end{tabular}
  \par\vspace{0.4ex}%
  \begin{minipage}{\linewidth}
    \centering\footnotesize
    $^\dagger$Proposed complexity--semantic router.
    Pooled EM: 48.1\%; pooled mean cost: 1.616\,mUSD/query ($N{=}965$).
  \end{minipage}
\end{table*}

\begin{table}[t]
  \centering
  \footnotesize
  \setlength{\tabcolsep}{2pt}
  \caption{Benchmark transfer: exact match (\%) on executed workloads (summary).}
  \label{tab:benchmark-transfer-em}
  \label{tab:deployment-em}
  \begin{tabular}{@{}lrrrr@{}}
    \toprule
    Workload & Proposed & Heur.\ baseline & ReAct & CoT \\
    \midrule
    GAIA & 18.8 & 24.9 & \textbf{26.1} & 6.1 \\
    HotpotQA & 58.5 & \textbf{64.0} & \textbf{69.0} & 36.5 \\
    MuSiQue & 36.0 & 35.0 & \textbf{36.5} & 2.5 \\
    MMLU & 55.0 & \textbf{56.0} & \textbf{59.0} & 53.0 \\
    MATH & 67.0 & 67.0 & 63.0 & \textbf{69.5} \\
    Pooled & 48.1 & \textbf{50.3} & {---} & {---} \\
    \bottomrule
  \end{tabular}
\end{table}

Table~\ref{tab:live-benchmark} reports exact match and mean cost across five benchmark
families; Table~\ref{tab:benchmark-transfer-em} summarizes exact match for the proposed
router, a heuristic baseline, and representative fixed strategies.
We distinguish in-domain workloads used in routing-benchmark construction (HotpotQA,
MuSiQue, MATH) from held-out families (GAIA, MMLU-Pro).
The proposed router attains 48.1\% pooled exact match versus 50.3\% for the heuristic
baseline; it does not uniformly outperform heuristic routing or the strongest fixed strategy
per family (Table~\ref{tab:benchmark-transfer-em}).
ReAct~\citep{yao2023react} remains the strongest fixed strategy on several workloads
(e.g., 69.0\% on HotpotQA).
\textbf{These results suggest} that routing-benchmark improvements do not fully translate into
exact-match gains across all benchmark families.
Distribution shift---especially on held-out GAIA and MMLU-Pro workloads not used in routing-benchmark
construction---and the gap between regret (which penalizes cost) and raw exact match (which does
not) both contribute to this pattern.
ReAct remains a strong \emph{fixed} policy on retrieval-heavy workloads, indicating that routers
must overcome both representation limits and supervision--metric mismatch to beat a competent
static default.

\subsection{Discussion}
\label{sec:discussion}

\paragraph{What works.}
On the routing benchmark, cost-soft supervision yields the largest reduction in utility regret
(agent regret $0.361\to 0.192$; executed exact match 63\%$\to$82\%; Table~\ref{tab:hard-vs-soft}).
The best configuration combines cost-soft training with the complexity--semantic representation
(total regret $0.198$; Tables~\ref{tab:hard-vs-soft} and~\ref{tab:feature-regret}).

\paragraph{What partially works.}
Planner-derived complexity provides complementary routing information when fused with semantic
embeddings, but not in isolation: complexity-only features attain higher regret than
embedding-only baselines ($0.273$ vs.\ $0.242$; Table~\ref{tab:feature-regret}).
Supervision design contributes more to regret reduction than feature enrichment alone, though
both are needed for the best observed configuration.

\paragraph{What remains difficult.}
Routing-benchmark improvements do not fully translate into exact-match gains across all
benchmark families (48.1\% pooled exact match vs.\ 50.3\% for a heuristic baseline;
Table~\ref{tab:benchmark-transfer-em}).
Pooled utility is also slightly lower on transfer workloads ($0.440$ vs.\ $0.458$;
Table~\ref{tab:pooled-utility} in Appendix~\ref{app:benchmarks}).
The framework is most reliable where training and evaluation share oracle multi-strategy
supervision; transfer to executed public workloads should be validated per target family.

Routing policy analysis, calibration diagrams, cost--accuracy frontiers, and embedding-dimension
sweeps appear in Appendix~\ref{app:additional-analyses}.
Extended interpretation is in Appendix~\ref{app:discussion}.
